\section{Finite-scale radial comparison and the complete positive remainder}
\label{sec:v69-finite-scale-coverage}
A comparison on an entire fixed radial interval need not use the
analytic germ at its center. We formulate it directly in terms of
the exact physical angular sets. This includes portions beyond a
certified germ disk and labels which are absent near the center but
appear farther out. The resulting finite-scale condition and its
weighted tail are kept distinct from the earlier persistence number.

Fix $0<S<h_\chi$ and an interval $I=(a,b)$ with
$0\le a<b\le S$. Put
\begin{align}
 L_{z,\lambda}(u)&=\ell_{z,\lambda}(\sqrt{2u}),\qquad
 A_{z,\lambda}(S)=\frac1S\int_0^S L_{z,\lambda}(u)\,du,
                                      \label{eq:v69-radial-average}\\
 U_{z,\lambda}(I)&=
          \operatorname*{ess\,sup}_{u\in I}L_{z,\lambda}(u),
 \qquad
 V_{z,\lambda}(I,S)=\frac{U_{z,\lambda}(I)}{A_{z,\lambda}(S)}.
                                      \label{eq:v69-radial-ratio}
\end{align}
Set $V=0$ if $A=0$. In that case $L=0$ almost everywhere on
$(0,S)$, hence on $I$, and there is no source on this radial interval.
These quantities are defined on the original physical indicator,
not on a Taylor polynomial. They have no guard denominator.

\begin{lemma}[Finite-scale geometric certificate]
\label{lem:v69-finite-scale-certificate}
At every fixed chart and exact label, $V(I,S)$ is finite and
nonnegative. If a verified computation of the physical angular sets
gives $U(I)\le u_+$ and $A(S)\ge a_->0$, then
$V(I,S)\le u_+/a_-$. On the common coarea representative,
\begin{equation}\label{eq:v69-finite-scale-comparison}
 \operatorname*{ess\,sup}_{u\in I}
 b_{z,\lambda}^{\varepsilon,\mathrm{clr}}(t_z+u)
 \le G_\chi(b,S)V_{z,\lambda}(I,S)
       \min\{1,g_z+L_g\sqrt{2b}\}\widehat q_{z,S;\lambda},
\end{equation}
where
$G_\chi(b,S)=\exp\{K_\chi(\sqrt{2b}+\sqrt{2S})\}$.
No analyticity or fixed radial ordering of the physical angular
boundaries on $(0,S)$ is assumed.
\end{lemma}
\begin{proof}
The measurable angular fraction lies in $[0,1]$. If its average is
positive the quotient is finite; if the average is zero the preceding
zero convention is justified by nonnegativity and Fubini's theorem.
By \eqref{eq:v69-unguarded-profile},
\[
 \widehat q_{z,S;\lambda}\ge
          e^{-K_\chi\sqrt{2S}}J_z A_{z,\lambda}(S),\qquad
 \operatorname*{ess\,sup}_{u\in I}\widehat b_{z,\lambda}(t_z+u)
       \le e^{K_\chi\sqrt{2b}}J_z U_{z,\lambda}(I).
\]
Divide when $A>0$ and apply
\eqref{eq:v69-absolute-guard-comparison}. The zero case has already
been settled. The certificate assertion is the same quotient
inequality. It specifies what must be verified on the physical
sets; no numerical certificate for all long Lorentz words is asserted.
\end{proof}

Let $b^{I,S,V\le T;\le\delta}$ be the original clearance source
restricted to those chart--label pairs with $V(I,S)\le T$ and
$g_z\le\delta$, and to roof offsets in $I$. The selection is made
at the fixed scale $S$ before the collision limit.

\begin{theorem}[Height on a complete finite radial interval]
\label{thm:v69-finite-scale-height}
For fixed admissible $\chi,\varepsilon,I,S,T$ and $0\le\delta\le1$,
\begin{equation}\label{eq:v69-finite-scale-height}
 \mathcal H(b^{I,S,V\le T;\le\delta})
 \le CTG_\chi(b,S)\min\{1,\delta+L_g\sqrt{2b}\}
                                      (b-a+S).
\end{equation}
In particular, for $I=(0,S)$ the entire selected radius
$0<F_z-t_z<S$ on this stratum satisfies
\begin{equation}\label{eq:v69-full-radius-height}
 \mathcal H(b^{(0,S),S,V\le T})
                    \le2CTS e^{2K_\chi\sqrt{2S}}.
\end{equation}
This includes source beyond an angular-germ disk, positive-radius
births of zero germs, and zero-center guards, when they belong to the
stated finite-scale stratum.
\end{theorem}
\begin{proof}
At roof $t$, the centers contributing to offsets in $(a,b)$ lie in
$[t-b,t-a]$, an interval of length $b-a$. Sum
\eqref{eq:v69-finite-scale-comparison} over these centers. Bound
$V$ by $T$ and the guard envelope by its displayed common bound.
Apply Theorem~\ref{thm:v69-comparison-collar} with collar width $S$
and interval length $b-a$. This proves the first statement. The
choice $a=0$, $b=S$, $\delta=1$ proves the second. The averaging
collar remains positive-width throughout. A fixed finite disjoint
partition into radial intervals may be treated in the same way and
its finitely many estimates added. No countable decomposition or
exchange of an infinite sum with a collision limsup is used.
\end{proof}

\subsection{A physically weighted moment sufficient for the radial tail}
For the whole radial interval write $V_z(S)=V_{z,\lambda}((0,S),S)$.
For $\alpha>0$ define the positive trace moment
\begin{equation}\label{eq:v69-radial-moment}
 \mathcal M_\alpha(S)=\limsup_{m\to\infty}\sup_{R,\lambda,t}m^2
 \sum_{z:t_z\in[t-S,t]} V_z(S)^{1+\alpha}
                                     \widehat q_{z,S;\lambda}.
\end{equation}
The sum uses the actual physical collar mass with the original exact
label, divided by $S$. It is not an unweighted word count or a sum
of full reversible coefficients. The moment is allowed to be infinite.

\begin{proposition}[Weighted radial-tail transfer]
\label{prop:v69-radial-tail-transfer}
The original clearance source on $0<F_z-t_z<S$ with $V_z(S)>T$
satisfies, for every $T>0$,
\begin{equation}\label{eq:v69-radial-tail-transfer}
 \mathcal H(b^{(0,S),S,V>T})
 \le e^{2K_\chi\sqrt{2S}}T^{-\alpha}\mathcal M_\alpha(S).
\end{equation}
Thus a finite bound on this physically weighted moment gives ordered
tightness of this radial cutoff. No such uniform moment bound is
asserted for the Lorentz family by the present proposition.
\end{proposition}
\begin{proof}
Use \eqref{eq:v69-finite-scale-comparison} with $I=(0,S)$ and
bound the guard envelope by one. On $V>T$,
$V\le T^{-\alpha}V^{1+\alpha}$. Sum at the prescribed roof, whose
contributing centers are in $[t-S,t]$, and take the stated ordered
limsup. The argument is valid in the extended nonnegative reals;
it asserts a finite tail estimate only when the displayed moment
is finite.
\end{proof}

\subsection{One disjoint source partition}
Fix $0<H<S<h_\chi$, $2H<h_\chi$, and $2K\sqrt H<1$.
First take the actual source $b^{\rm a,K,H}$ of
\eqref{eq:v69-small-center-source}. Among all remaining physical
points in the selected charts with $0<F_z-t_z<S$, take those whose
chart--label pair has $V_z(S)\le T$ and call their source $o_{69}$.
Call the rest within that radius $t_{69}$. Finally let $e_{69}$
contain every point not yet assigned: the original outside source
and all portions of the selected charts with $F_z-t_z\ge S$.
Then, as densities of disjoint positive restrictions,
\begin{equation}\label{eq:v69-complete-positive-partition}
 b^{\varepsilon,\mathrm{clr}}
            =b^{\rm a,K,H}+o_{69}+t_{69}+e_{69}.
\end{equation}
The decomposition preserves the first-clearance witness, the
next-collision mark, the original occupation and terminal convention,
and every arithmetic zero class. In particular zero-center source
is not discarded. Outside an inner angular certificate it is assigned
by the finite-scale partition just described.

The first two restrictions have unconditional bounds
\[
 \mathcal H(b^{\rm a,K,H})\le6CH F_{\chi,K}(H),\qquad
 \mathcal H(o_{69})\le2CTS e^{2K_\chi\sqrt{2S}}.
\]
The second follows by positive domination by the full $V\le T$
source, even though an overlap was removed. Similarly $t_{69}$ is
dominated by the full $V>T$ source in
Proposition~\ref{prop:v69-radial-tail-transfer}. Consequently the
unchanged pointwise error satisfies, for every fixed reconstruction
band $B$ and its width $\varepsilon(B)=A_0B^{-1/12}$,
\begin{align}
 \mathcal E_M\le{}& C_MB^{-1/192}
       +\mathcal H(b^{\varepsilon(B),\mathrm{inc}})
       +6CH F_{\chi,K}(H)\notag\\
       &+2CTS e^{2K_\chi\sqrt{2S}}
         +\mathcal H(t_{69})+\mathcal H(e_{69}).
                                      \label{eq:v69-canonical-endpoint-budget}
\end{align}
A proved bound on $\mathcal M_\alpha(S)$ may replace the fifth
term by \eqref{eq:v69-radial-tail-transfer}; it may not be silently
assumed. First fix $B,\chi,K,H,S,T$, then take the collision limsup.
Any subsequent choices may depend on $B$, but not on $m$.

The finite-scale theorem controls the entire chosen radial interval
on its stated stratum, rather than treating its outer part as
controlled by a shrinking inner collar. It does not estimate the
high-$V$ tail without the moment, nor the exterior $e_{69}$.
The latter still contains failed taper, selected grazing, noncritical
rank, uncovered states, and all omitted outer annuli. The complete
first-incidence height remains a separate positive term. Neither
its exponential finite-count estimate in
Section~\ref{sec:v61-inverse-incidence}, nor fixed-count finiteness
of $V$, supplies the missing ordered estimate.

The original unrestricted two-sided pointwise theorem and its
same-roof bridge, essential-likelihood and conditioned-path
consequences retain precisely these complete-source obligations.
The new route is an enlargement of the controlled Lorentz source,
not a replacement of that endpoint by a different model or topic.
\label{end:v69-new-mathematics}
